\section{Full generator equations for all eleven families}
\label{sec:all-generators}
Section~\ref{sec:array-repr} gave the two most-used families (square, hexagonal) in
full. This section derives all eleven implemented families
(\texttt{array\_calculation/generators/}, verified read directly file-by-file; mtimes
in Chapter~\ref{ch:provenance}), since Chapters~\ref{ch:model-a}--\ref{ch:model-b}
evaluate every one of them and Section~\ref{sec:hex-intuition}'s symmetry argument only
has force once the alternatives are actually written down.

Every generator shares one post-processing step, \texttt{base\_generator.finalize()}
(verified read directly): greedy overlap rejection (points are kept in generation
order, any later point closer than $a_{\mathrm{hole}}$ to an already-kept point is
discarded), then recentering on the surviving set's own centroid. This shared step is
the reason a truncated (non-lattice-complete) candidate set can end up off-centre --
the mechanism behind Section~\ref{sec:centered-bug} below -- and is stated once here
rather than repeated in every family's own derivation.

\subsection{Square}
\begin{equation}
(x_{m,n}) = R(\theta_0)\cdot p\,(m,n), \qquad (m,n)\in\mathbb{Z}^2
\end{equation}
Free parameter: pitch $p$ (nearest-neighbor spacing), plus optional rotation
$\theta_0$. Construction: a square meshgrid large enough to contain $N$ points,
radially sorted by distance from the origin, the closest $N$ kept, then rotated by
$\theta_0$ (rotation is applied \emph{after} the radial sort, so it changes only each
point's $(x,y)$, never which points are selected -- rotation preserves every pairwise
and every origin distance exactly).

\subsection{Hexagonal}
\begin{align}
y_r &= r\,p\,\frac{\sqrt3}{2}, \qquad r\in\mathbb{Z} \\
x_{r,m} &= mp + \begin{cases}p/2 & r\text{ odd}\\0&r\text{ even}\end{cases}, \qquad m\in\mathbb{Z}
\end{align}
Row height $p\sqrt3/2$ and alternate-row half-pitch offset are the standard hexagonal
close-packing construction; identical radial-sort-then-keep-closest-$N$ selection as
square.

\subsection{Rectangular}
\begin{equation}
(x_{m,n}) = R(\theta_0)\cdot(mp_x, np_y), \qquad (m,n)\in\mathbb{Z}^2
\end{equation}
Independent $x$- and $y$-pitches $p_x,p_y$. At $p_x=p_y$ this reduces to \emph{exactly}
the same lattice as square's own construction -- the two generators use deliberately
structurally identical code (same meshgrid-then-radial-sort approach) specifically so
this equivalence holds exactly, not approximately, and is checked directly in
\texttt{spot\_check\_generators.py}.

\subsection{Concentric}
\label{sec:concentric-family}
Ring $k=0$ is a single hole at the origin. Ring $k\ge1$ sits at radius $r_k=k\,s$ (ring
spacing $s$) with
\begin{equation}
n_k = \mathrm{round}(n_0\,k), \qquad
\theta_{k,j} = \frac{2\pi j}{n_k},\ j=0,\dots,n_k-1
\end{equation}
hole count \emph{scaling with ring circumference} (constant angular spacing between
adjacent holes on every ring, equal to ring 1's own spacing $2\pi/n_0$) rather than a
fixed count per ring -- a fixed count would make outer rings implausibly sparse or
force inner-ring holes closer than $a_{\mathrm{hole}}$ (which \texttt{finalize()}'s
overlap rejection would then discard anyway), so the constant-angular-spacing rule is
the more physically sensible default for a fabricatable pattern (source: module
docstring, this exact reasoning, verified read directly).

\subsection{Ring-offset concentric}
\label{sec:ring-offset-family}
Identical ring radius/count formula to Section~\ref{sec:concentric-family}, with one
addition: odd-indexed rings ($k$ odd) receive a fixed half-angular-step rotation
$\pi/n_k$ relative to their unrotated position; even rings receive none --
\begin{equation}
\theta_{k,j} = \frac{2\pi j}{n_k} + \begin{cases}\pi/n_k & k\text{ odd}\\0&k\text{ even}\end{cases}
\end{equation}
the ``brick-course'' pattern: adjacent rings' holes sit staggered rather than radially
aligned. An \emph{alternating} offset (not a progressively-accumulating one) was
chosen as the more literal reading of a brick pattern -- real brick courses alternate,
they do not progressively drift -- and because it is the simpler, unambiguous
construction; a progressive variant was considered and explicitly not implemented,
per the module's own stated reasoning against shipping an unused second flag.

\subsection{Concentric-centered, and a real, documented bug it exists to fix}
\label{sec:centered-bug}
Investigating this family's own motivation directly (rather than assuming it), its
author found that \texttt{concentric.py} and \texttt{ring\_offset\_concentric.py}
\emph{already} start their own point list with a literal centre hole at $(0,0)$ -- so
the reason neither family's \emph{final}, post-\texttt{finalize()} output has a hole at
the origin is not a missing centre point, but \texttt{base\_generator}'s own shared
recentering step: a \emph{complete} ring sums to exactly $(0,0)$ by rotational
symmetry and never disturbs this, but a \emph{truncated} final ring (the common case,
whenever the requested $N$ does not exactly fill an integer number of complete rings)
does not sum to zero, and the shared recenter-on-centroid step then shifts the nominal
centre hole away from the true origin by exactly that leftover imbalance.

\begin{warnbox}{Confirmed against a real, already-fabrication-relevant case}
\texttt{ring\_offset\_concentric}'s own $N=20$, ring spacing $4000\,\mu$m, $n_0=8$
candidate -- its real $M=2$--$3$D winner in an earlier evaluation pass -- has its
nominal centre hole recentred to \textbf{$1704.8\,\mu$m} from the true origin.
Reproduced bit-for-bit by a naive ``prepend then recenter together'' re-implementation,
confirming this mechanism directly, not by inference.
\end{warnbox}

\textbf{Fix, and the one substantive difference this generator adds}: recenter the
ring points on their \emph{own} centroid first, \emph{then} prepend the $(0,0)$ centre
hole, \emph{then} run \texttt{finalize()} once more on the combined set. Since the ring
points are already exactly zero-mean before the centre hole is added, and the centre
hole itself contributes $(0,0)$ to any mean, the combined set's mean is exactly
$(0,0)$ by construction -- the final recentering step becomes provably a no-op. Checked
directly on the identical displaced case above: centre-hole displacement drops to
$1.6\times10^{-13}\,\mu$m (floating-point noise), confirming the fix, not merely
asserting it.

\subsection{Archimedean spiral}
$r=c\theta$ (unwrapped, not reduced mod $2\pi$) -- the curve where each full
revolution increases $r$ by the constant amount $2\pi c$. Points are placed along the
curve at approximately constant arc-length spacing: since $ds/d\theta=\sqrt{r^2+c^2}$
for this curve, a first-order (Euler) angular step
\begin{equation}
d\theta = \frac{a_{\mathrm{hole}}\cdot f_{\mathrm{spacing}}}{\sqrt{r^2+c^2}}
\end{equation}
(evaluated at the current $r$ before each step, $f_{\mathrm{spacing}}=1.5$ by default)
keeps successive points near-evenly spaced along the arc. Explicitly \emph{not} the
same construction as the Fibonacci/golden-angle family below -- this generator advances
along one continuous outward-spiralling arm; the golden-angle family jumps by a fixed
large irrational angle each step and never traces a continuous curve at all. Confirmed
different point clouds directly, not merely by differing formulas
(\texttt{spot\_check\_generators.py}'s archimedean-vs-fibonacci difference check).

\subsection{Fibonacci (Vogel) spiral and golden-angle phyllotaxis}
\begin{equation}
r_n = c\sqrt n, \qquad \theta_n = n\,\alpha_g, \qquad \alpha_g = \pi(3-\sqrt5)\approx137.50776^\circ,
\qquad n=1,\dots,N
\end{equation}
$\alpha_g=2\pi(1-1/\varphi)$, $\varphi$ the golden ratio -- the irrational angle that
maximizes packing uniformity for this family (Vogel, 1979). ``Golden-angle
phyllotaxis'' and the ``Fibonacci/Vogel spiral'' are, mathematically, the
\emph{identical} formula -- there is no separate standard point-packing construction in
the literature distinct from this one (the visually different ``nested-squares
Fibonacci spiral'' is an unrelated aesthetic curve, not a packing formula, and is not
what either generator name refers to here). Rather than risk two copies silently
drifting apart under future edits, \texttt{golden\_angle.py} re-exports
\texttt{fibonacci\_spiral.generate()} directly under its own family label; for the same
$(N,a_{\mathrm{hole}},c)$ the two generators' output coordinates are, by construction,
identical.

\subsection{Poisson-disk (blue-noise)}
Stochastic; minimum pairwise distance $r=a_{\mathrm{hole}}f_{\mathrm{spacing}}$
enforced via Bridson's algorithm (Bridson, SIGGRAPH 2007 sketches) rather than naive
dart-throwing-with-rejection, whose rejection rate is unbounded as the domain fills
(each new random dart becomes increasingly likely to land too close to an existing
point). Bridson's method instead grows the set from an active frontier -- only ever
proposing new candidates in the annulus $[r,2r]$ around an existing active point,
checked against neighbours via a background grid of cell size $r/\sqrt2$ (small enough
that each cell holds at most one sample) -- and terminates deterministically once no
active point can produce a valid candidate after $k=30$ tries, giving $O(1)$ amortized
work per accepted or retired point rather than the unbounded degradation of naive
rejection sampling. The domain is filled to near-maximal packing, then the $N$ points
closest to the origin are kept (same selection convention as square/hexagonal). Because
the algorithm is stochastic, \texttt{seed} is a required, not optional, parameter for
every call to be reproducible.

\subsection{Jittered lattice}
Starts from \texttt{hexagonal.generate()}'s own lattice, then perturbs every hole by a
random offset sampled \emph{uniformly by area} within a disk of radius
$f_{\mathrm{jitter}}\,p$:
\begin{equation}
\rho = f_{\mathrm{jitter}}\,p\,\sqrt{u}, \qquad u\sim\mathrm{Uniform}(0,1), \qquad
\phi\sim\mathrm{Uniform}(0,2\pi)
\end{equation}
The $\sqrt{u}$ is load-bearing, not optional: sampling the offset magnitude directly as
$f_{\mathrm{jitter}}\,p\,u$ would over-concentrate perturbed points near zero offset,
since area scales as $r^2$, not $r$ -- the same disk-sampling construction already used
elsewhere in this project's own Monte Carlo validation work
(\texttt{array\_optimization/vignetting.py}). A \emph{uniform}, not Gaussian, offset
distribution was chosen specifically because the parameter's own contract is ``offset
up to $f_{\mathrm{jitter}}\,p$'' -- a hard, stated bound; a Gaussian has unbounded
support, which would make that ``up to'' contract false in the tail, however unlikely.
The jittered set is then re-run through \texttt{finalize()}'s own overlap rejection
(jitter can push two holes that were exactly at the lattice's own valid spacing closer
than $a_{\mathrm{hole}}$ -- these must be rejected exactly as any other overlapping
candidate would be, not silently kept because they originated from a valid lattice).
At $f_{\mathrm{jitter}}=0$ every offset is identically zero, so this reduces exactly to
\texttt{hexagonal.generate()}'s own output -- checked directly, not merely argued, in
\texttt{spot\_check\_generators.py}.
